\section*{Заключение}
\addcontentsline{toc}{section}{Заключение}

В рамках лабораторных работ~1--5 реализован полный цикл работы
с графом: генерация связного ациклического графа по распределению
Рэлея--Райса с softmax-нормировкой степеней, вычисление характеристик
(эксцентриситеты, радиус, диаметр, центр), построение матриц
минимальных и максимальных маршрутов длины~$k$ методом Шимбелла,
перечисление всех простых путей между двумя вершинами, обход рёбер
графа в ширину, нахождение кратчайшего пути алгоритмом Дейкстры
с выводом вектора расстояний и сравнение алгоритмов по числу итераций.
Исходный код реализации всех лабораторных работ размещён в открытом
github-репозитории~\cite{github}.

В лабораторной работе~3 на основе порождённого графа формируется
транспортная сеть: строятся матрицы пропускных способностей и
стоимостей дуг, находится максимальный поток и поток минимальной стоимости величиной
$\lfloor\tfrac{2}{3} \cdot F_{\max}\rfloor$ методом последовательных
кратчайших путей на базе алгоритма Беллмана--Форда.

В лабораторной работе~4 для неориентированного графа вычисляется
число остовных деревьев по матричной теореме Кирхгофа, строится
минимальный остов алгоритмом Борувки, полученное дерево кодируется
и декодируется кодом Прюфера с сохранением весов рёбер, а также
находится минимальная раскраска графа методом перебора с возвратом.

В лабораторной работе~5 реализована проверка графа на эйлеровость с автоматической достройкой графа и нахождением эйлерова цикла алгоритмом Иерхольцера. Также построена фундаментальная система разрезов (ФСР) на базе остова Борувки с вычислением симметрической разности для любого выбранного подмножества разрезов.

Все пункты заданий выполнены. Граф из лабораторной~1
переиспользуется в последующих работах без перегенерации. Результаты
каждого шага сохраняются в текстовые файлы и передаются
в Python-скрипты для визуализации без перезапуска исполняемого файла.

\subsection*{Сильные стороны реализации}

Архитектура разделена на независимые слои: граф не знает о
генераторе, генератор не знает об алгоритмах, алгоритмы не
знают о визуализации. Это позволяет заменять любой компонент
без изменения остальных.

Шаблонный \texttt{GraphBase} даёт одну кодовую базу для графа
и сети потоков за счёт статического полиморфизма. Шаблонный
генератор с лямбдами переиспользуется для обоих типов графов
и любого закона генерации весов без дублирования алгоритма построения.

\texttt{std::optional} в матрицах Шимбелла и в интерфейсе графа
исключает специальные значения и делает семантику «ребра нет» и
«пути нет» явной на уровне типов.

Наличие веб-интерфейса с полной визуализацией и анимацией алгоритмов
на графах.

\subsection*{Слабые стороны реализации}

Метод Шимбелла работает за $O(k \cdot n^3)$ и хранит две матрицы
$n \times n$ в памяти.

Поиск всех простых путей имеет факториальную сложность. На разреженных
ациклических графах это допустимо, но никакой защиты при больших
входных данных не предусмотрено.

Минимальная раскраска реализована полным перебором с возвратом
со сложностью $O(n \cdot k^n)$ в худшем случае~--- алгоритм
применим лишь на небольших графах и не масштабируется на плотные
графы с большим хроматическим числом.

Передача данных между C++ и Python реализована через текстовые файлы:
C++ записывает результат в \texttt{assets/txt/}, Python читает его
при запуске скрипта. Формат намеренно простой для учебной задачи.
В production-окружении такое взаимодействие между процессами
заменяется брокерами сообщений (Kafka, RabbitMQ) или очередями задач
(Celery, ZeroMQ)~--- они обеспечивают типизированные сообщения,
подтверждение доставки и горизонтальное масштабирование.

\subsection*{Масштабирование}

Замена файлового обмена на брокер сообщений развяжет процессы по
времени, позволит запускать несколько воркеров визуализации параллельно
и добавит гарантии доставки.

Переход на WebSocket позволит C++ передавать результат напрямую
в браузер: визуализация появится сразу по завершении вычислений
без перезагрузки страницы.

Замена статичных PNG на интерактивные графы через Plotly~--- с
возможностью масштабирования, выделения вершин и фильтрации рёбер
по весу прямо в браузере.

Параллелизация алгоритмической части: метод Шимбелла считает $n$
независимых строк матрицы на каждом шаге, поиск всех путей разбивается по стартовым рёбрам, а итерации алгоритма.

Покрытие тестами GTest граничных случаев генератора, корректности
матриц Шимбелла, полноты поиска путей, инвариантности веса при
кодировании и декодировании кода Прюфера и правильности раскраски.
